\citet{biasi_what_2025} ask whether and for whom school capital investments financed through local bond elections improve student achievement and are valued by households. Methodologically, they assemble the first near-national panel of bond referenda linked to district finances, test scores, and house prices, and they identify dynamic effects with a stacked regression discontinuity design built around state-specific approval thresholds. Substantively, they find that bond passage raises capital spending, improves test scores over the medium run, and is capitalized into house prices on average, with returns that differ by project type and district characteristics. This finding matters because school capital is a large, locally financed, and recurrent expenditure whose returns had previously been estimated mainly within single states \citep{cellini2010value,hong2016investing,rauscher2020delayed}; a credible national estimate is directly relevant to education-finance policy. This replication reproduces the central first-stage and reduced-form estimates, then tests their robustness to local close-election comparisons and their external validity across states.

\subsection{Data}

The analysis relies on the national panel of school capital elections and outcomes constructed by \citet{biasi_what_2025}. The original authors assemble a database of school bond referenda from administrative records and public information requests, covering local bond elections in 41 U.S. states between 1995 and 2017. After the sample restrictions used in their analysis, the main replication sample contains 29 states.

The bond election data are linked to annual district-level finances from the NCES Annual Survey of School System Finances and the Census of Governments, student achievement data from the Stanford Education Data Archive, and district-level housing price indices from \citet{contat_larson_2024}. The resulting panel allows the original paper to estimate fiscal, academic, and housing-market responses to capital projects authorized through local bond elections.

% This replication uses the cleaned master files in the \citet{biasi_what_2025} replication package. 
I do not alter the original data construction for the computational replication. For the state-level extension, I use the same stacked files but estimate the main specifications separately by state after applying the same McCrary-based state screen and the additional state-level sample restrictions documented in Appendix Table \ref{tab:state_sample_audit}.

\subsection{Stacked Dynamic RD Framework}

The original study exploits discontinuities in bond authorization generated by state-specific approval thresholds. Let $m_{jc}$ denote the vote margin in district $j$ and election cohort $c$, defined as the difference between the yes-vote share and the required approval threshold. A bond is authorized when $m_{jc} \geq 0$. Outcomes $Y_{jct}$ include capital expenditures per pupil, standardized test scores, and a district-level house price index.

Because districts can hold multiple elections over time, \citet{biasi_what_2025} use a stacked dynamic regression discontinuity design. Each election is treated as a cohort, and districts are followed in event time around that election. The stacked panel is indexed by district, cohort, and calendar year. The estimating equation includes district-by-cohort fixed effects, state-by-year fixed effects, event-time treatment indicators, flexible functions of the vote margin, and controls for other bond elections in the district's history.

In compact form, the main OLS specification is
\[
Y_{jct}
= \alpha_{jc} + \gamma_{s(j)t}
+ \sum_{\tau \neq 0} \beta_{\tau}
\mathbf{1}\{t-c=\tau\}\mathbf{1}\{m_{jc}\geq 0\}
+ X_{jct}'\delta + g_c(m_{jc}) + u_{jct},
\]
where $\alpha_{jc}$ are district-by-cohort fixed effects, $\gamma_{s(j)t}$ are state-by-year fixed effects, $X_{jct}$ contains bond-history controls, and $g_c(m_{jc})$ denotes cohort-specific vote-margin polynomials. Standard errors are clustered at the district level.

Two features of this estimand are worth making explicit, because they motivate the robustness audit in Section \ref{sec:bandwidth}. First, the $\beta_\tau$ identify a treatment effect at the threshold under the standard RD continuity and no-manipulation conditions \citep{lee2008randomized,lee_lemieux2010}, whose credibility in close-election settings is supported by the smoothness of the running-variable density \citep{mccrary2008manipulation,cattaneo2020simple} and by the broad empirical validity of close-election designs documented in \citet{eggers2015validity}. Second, because $g_c(\cdot)$ is fit on the full sample, the local effect is recovered in part by extrapolating the vote-margin polynomials toward the cutoff rather than by comparison among near-threshold elections alone---the efficiency-versus-locality trade-off emphasized by \citet{calonico2014robust} and \citet{imbens2012optimal}. The full-sample specification buys precision at the cost of leaning on functional form, which is exactly the assumption a margin-restriction audit probes.

% \subsection{Replication and Extensions}

% The three exercises map onto the replication categories recognized by \emph{Economics of Education Review} and formalized by \citet{clemens2017meaning}. The computational replication is a \emph{verification}: I implement the original Stata workflow to regenerate the main first-stage and reduced-form estimates and compare them with an independent R implementation of the central tables, checking that the published numbers can be reproduced from the package.

% The vote-margin robustness exercise keeps the BLS full-sample specification as the baseline and re-estimates the same stacked DRD model under symmetric restrictions on the absolute vote margin: 20, 15, 10, and 5 percentage points. These estimates are reported as robustness checks rather than replacements for the published specification.

% The state-level extension estimates the same dynamic treatment effects separately by state for states with sufficient usable data. The state-specific estimates are then linked to descriptive state-level averages of project categories, baseline capital stock, free or reduced-price lunch share, and minority share. These mechanism exercises are descriptive because state-level characteristics are not experimentally assigned and the number of states is limited.
